\section{A full topology kernel from quadrilateral content}
\label{sec:kernel}

The direct two-weight census still carries the numerical size of $x$ through
its orbit queries. The existing disc-count kernel removes this cost when
large vertex links and a common quadrilateral factor conceal a small core.
The type spectrum allows that reduction to be extended to the entire
abstract topology, without the incorrect rule that every component count
scales linearly.

\subsection{Peeling vertex links and proving divisibility}

For each global triangulation vertex $v$, let
\[
 m_v=\min\{x_{t,a}:\text{tetrahedron corner }(t,a)\text{ belongs to }v\}.
\]
Here $x_{t,a}$ denotes a triangle coordinate. Let $L_v$ be the vertex-link
vector, with one triangle at each such corner. Form
\[
 y=x-\sum_v m_v L_v.
\]
The represented links can be placed in disjoint small vertex neighbourhoods,
disjoint from the remaining surface, so this subtraction has a geometric
disjoint-union interpretation. Each $L_v$ is a disc at a boundary vertex and
a sphere at an interior vertex, by the validated manifold-link conditions.
All matching equations and quadrilateral constraints remain satisfied.
For each global vertex, at least one incident triangle coordinate of $y$
is zero.

Let $q_1,\ldots,q_{3t}$ be the quadrilateral coordinates, and set
\[
 g=\gcd(q_1,\ldots,q_{3t}),\qquad Q=\max_i q_i.
\]
We take $g=0$ when all quadrilateral coordinates vanish.

\begin{lemma}[Quadrilateral core]
\label{lem:quad-core}
If $g=0$, then $y=0$. If $g>0$, every coordinate of $y$ is divisible by
$g$ and the vector $p=y/g$ is integral and admissible. Moreover,
\begin{equation}
 \max_i p_i\le (4t-1)Q/g.
 \label{eq:core-bound}
\end{equation}
\end{lemma}
\begin{proof}
Across a glued face, the matching equation for an arc cutting off the
identified corner has the form
\[
 T_1+Q_1=T_2+Q_2,
\]
where each $T_i$ is the relevant triangle count and each $Q_i$ is the
quadrilateral count that contributes to that arc. Consequently
$T_1-T_2=Q_2-Q_1$. The corner-adjacency graph of a global vertex is
connected because its link is a connected surface. It has at most $4t$
vertices. After peeling, one of its triangle counts is zero.

If all quadrilaterals vanish, the triangle count is constant along every
edge of this graph, so all its triangle counts are zero. This proves the
first assertion. Otherwise, all differences are divisible by $g$.
Propagating from a zero triangle proves divisibility of every triangle;
quadrilateral divisibility is the definition of $g$.

Each $Q_i$ lies between zero and $Q$, so an edge difference has absolute
value at most $Q$. A simple path from the zero triangle has at most
$4t-1$ edges. Thus every peeled triangle is at most $(4t-1)Q$.
The same upper bound covers the quadrilaterals. Division by $g$ gives
\eqref{eq:core-bound}. Homogeneity preserves matching, nonnegativity,
and the quadrilateral constraints.
\end{proof}

The arithmetic reduction and this core bound are inherited from the
repository's report 53 integration \cite{proveit-baseline}; we include the
proof to state exactly what the full-spectrum extension relies upon.

\subsection{Normal scaling of a disconnected surface}

\begin{lemma}[Ordered-sheet scaling]
\label{lem:scaling}
In an orientable ambient three-manifold, $g$ times the normal vector of an
orientable connected surface contributes $g$ copies of that surface.
For a nonorientable connected surface it contributes
$\lfloor g/2\rfloor$ copies of its connected orientation double, and one
copy of the original surface if $g$ is odd.
\end{lemma}
\begin{proof}
Over each normal disc, the $g$ local sheets have an order
$j=0,\ldots,g-1$. Across any face, the order is preserved or reversed.
Normal fibre monodromy is therefore generated by
$j\mapsto g-1-j$. For a two-sided component there is no nontrivial
monodromy, so each sheet closes separately. For a one-sided component,
the orbits of this reversal are the pairs $\{j,g-1-j\}$, together with
a central singleton when $g$ is odd. Each pair is the connected normal
orientation double. Orientability and two-sidedness coincide in the
orientable ambient manifold. Disjoint interval neighbourhoods of the
original components justify applying this argument componentwise.
\end{proof}

For a nonorientable component with crosscap number $k$ and $b$ boundary
circles, its double has signature $(2\chi,2b)$ and orientable genus $k-1$.
Twice a M\"obius band is thus one annulus; twice a Klein bottle is one
torus; twice a projective plane is one sphere.

\begin{theorem}[Complete spectrum restoration]
\label{thm:core}
Suppose $g>0$ and let $O_p(v),A_p(v)$ be the orientable and nonorientable
type counts of the core $p$ at signature $v=(\chi,b)$. Let
$d=\sum_{v\text{ boundary}}m_v$ and
$s_0=\sum_{v\text{ interior}}m_v$. Then the original spectrum is
\begin{align}
 O_x(v)={}&gO_p(v)+\lfloor g/2\rfloor A_p(v/2)
                +d\mathbf1_{v=(1,1)}+s_0\mathbf1_{v=(2,0)},
 \label{eq:restore-O}\\
 A_x(v)={}&(g\bmod2)A_p(v).
 \label{eq:restore-A}
\end{align}
An absent or nonintegral half contributes zero. If $g=0$, only the two
vertex-link terms remain. For $r$ core type rows, at most $2r+2$ output
type rows are required.
\end{theorem}
\begin{proof}
Peeling the vertex links gives a geometric disjoint union of those links
and $F(gp)$. Apply \cref{lem:scaling} to each connected component of
$F(p)$. An orientable component stays at its original signature; each
nonorientable orientation double appears at twice its original signature.
These contributions are exactly \cref{eq:restore-O,eq:restore-A}.
The final two terms add the peeled discs and spheres. Each orientable
core row contributes at most one type and each nonorientable row at most
two; the stated support bound follows after coalescing identical types.
\end{proof}

The torus--Klein zero-signature case is automatically correct in
\eqref{eq:restore-O}: a zero-signature nonorientable row contributes its
double at the same numerical signature, but with orientability changed.
The type key therefore includes orientability during restoration.

As a consistency check, total Euler characteristic and total boundary
circle count of the core scale by $g$. For a nonorientable component,
$2\lfloor g/2\rfloor+(g\bmod2)=g$. The \emph{number of components}
instead becomes
\[
 C(gp)=g\,C_{\rm or}(p)+\lceil g/2\rceil C_{\rm nonor}(p).
\]
It cannot in general be obtained by multiplying $C(p)$ by $g$.

\subsection{A family missed by whole-vector gcd reduction}

Let $D$ be a fixed vertex-reduced primitive normal meridian in a one-vertex
layered solid torus, and let $L$ be its vertex-link disc. Consider
\begin{equation}
 x_g=gD+(g+1)L.\label{eq:disc-family}
\end{equation}
The layered fixture has a quadrilateral coordinate one and a zero triangle
at the vertex. Thus $x_g$ contains coordinates $g$ and $g+1$, and its full
coordinate gcd is one. Whole-vector scaling finds no reduction. Peeling
the vertex link removes $g+1$ copies of $L$, and quadrilateral content is
$g$, so the query core is exactly $D$ for every $g$.

The new type spectrum has $2g+1$ disc components, represented by one row.
This row merges the $g$ essential meridians with the $g+1$ inessential
vertex-link discs. That is the correct abstract topology and also an
explicit demonstration that this output alone is not an essential-disc
certificate.

Replacing $D$ by a fixed vertex-reduced M\"obius band $M$ gives
\[
 gM+(g+1)L:
 \quad (g+1)\text{ discs},\quad
 \lfloor g/2\rfloor\text{ annuli},\quad
 (g\bmod2)\text{ M\"obius bands}.
\]
For a fixed closed Klein bottle $K$ and its compatible boundary vertex link,
the analogous family gives $(g+1)$ discs,
$\lfloor g/2\rfloor$ tori, and a Klein bottle precisely for odd $g$.
All these outputs retain constant type support while their multiplicities
grow through arbitrary binary integers.

\subsection{Bit complexity and exact savings}

Let $B_0=\max(1,\max_i\bits(p_i))$ for a nonzero core. From
\eqref{eq:core-bound},
\[
 B_0=O\bigl(\log(t+1)+\log(1+Q/g)\bigr).
\]
Let $V(t,B)$ denote the inherited source-validation and decoding cost,
and let $P_2(t,b)$ denote the polynomial cost of the direct two-weight
census on $b$-bit coordinates, including boundary markers and inversion.
The full core algorithm has the decomposition
\begin{equation}
 V(t,B)+\poly(t,B)+P_2(t,B_0)
 +O\bigl(r\M(B+\ell)+r(B+\ell)\log(r+1)\bigr),
 \label{eq:core-cost}
\end{equation}
where $\ell$ bounds the core spectrum fields. The additional polynomial
front-end term covers minima, gcd, exact division, and core matching
checks. The last term bounds scaling and merging the type rows.

For a fixed core and unbounded $g$ and peeled-link multiplicities, the
orbit-discovery schedules and their endpoints depend only on that fixed
core. The expensive local transcripts are independent of the outer
multiplicities. Input validation, fingerprints, gcd/division, final
multiplications, and serialization still depend on the original input
bits. In particular, \eqref{eq:core-cost} is not constant-time processing
of an arbitrarily long input. A certificate records the original source
digest, all peeled multiplicities, the exact divisor and core, the core
proof, and the restored spectrum. The outer fields continue to have the
necessary $O(tB)$-scale encoded contribution.

On primitive vectors with little or no peelable triangle mass, $B_0$ need
not be smaller than $B$. The core calculation then adds some overhead.
This parameter dependence is measured below and is part of the algorithm's
contract, not a hidden exception.

\section{What this changes in the global complexity analysis}
\label{sec:global}

The practical improvement removes two forms of unnecessary information
from a topology query: all $7t$ component coordinates when only abstract
type is requested, and repeated large multiplicity data inside every orbit
event when a small primitive core suffices. Neither operation reduces the
number of normal vectors an outer search may have to inspect.

\begin{proposition}[Composition with a controlled outer algorithm]
If an outer algorithm makes at most $Q(n)$ supplied-vector topology
queries, each of encoded input size at most $S(n)$, the present observers
contribute at most $Q(n)\poly(S(n))$ time, including validation and output.
In particular, quasi-polynomial bounds on both $Q(n)$ and $S(n)$ preserve
quasi-polynomial cost for this observer stage.
\end{proposition}
\begin{proof}
Apply \cref{thm:census} or \eqref{eq:core-cost} to each call and add their
costs. Products of a fixed number of functions of the form
$2^{(\log n)^{O(1)}}$ remain of that form.
\end{proof}

This proposition is a compositional guarantee, not a proof of its
hypotheses for unknot recognition. Three obligations remain especially
concrete. First, construct and certify the correspondence between an
input knot diagram and the particular triangulated exterior being
queried. Second, discover useful normal vectors or multi-surfaces without
an uncontrolled exponential search. Third, perform the necessary cutting,
retriangulation, compression, and hierarchy updates while controlling
their encoded sizes and the number of restarts. The present type spectrum
does not contain enough embedding or attachment data to discharge the
third obligation by itself.

The distinction also rules out a tempting but invalid recognizer: a
negative spectrum query on one vector does not show that no compressing
disc exists in the manifold. Even a positive disc row must be paired with
a source-bound boundary-essentiality argument before it becomes a
compressing-disc certificate. The existing torus-boundary parity observer
can provide that additional componentwise information for an appropriate
query; it cannot infer it from an unlabelled disc multiplicity alone.
